\subsection{Polynomials}

\begin{itemize}
    \item \textbf{Definition: } A \textbf{polynomial} is an expression of the form
          \[
              a_n x^n + a_{n-1} x^{n-1} + \cdots + a_1 x + a_0,
          \]
          where $a_n, a_{n-1}, \ldots, a_0$ are coefficients and $n$ is a non-negative integer called the degree of the polynomial.

    \item \textbf{Root of a polynomial: } A \textbf{root} of a polynomial $P(x)$ is a value $r$ such that $P(r) = 0$.

    \item \textbf{Fundamental Theorem of Algebra: } Every non-constant polynomial with complex coefficients has at least one complex root.

    \item \textbf{Factorization: } A polynomial of degree $n$ can be factored into linear factors over the complex numbers, i.e., if $P(x)$ is a polynomial of degree $n$, then it can be expressed as
          \[
              P(x) = a_n (x - r_1)(x - r_2) \cdots (x - r_n),
          \]
          where $r_1, r_2, \ldots, r_n$ are the roots of the polynomial (counting multiplicities).

    \item \textbf{Inversion: } Given a polynomial $P(x)$, the \textbf{inverse} of $P(x)$, denoted as $P^{-1}(x)$, is a polynomial such that $P(P^{-1}(x)) = x$ for all $x$ in the domain of $P^{-1}$.
\end{itemize}

\subsection{Algorithms}
\begin{itemize}
    \item \textbf{Fast Fourrier Transform (FFT): } An efficient algorithm to compute the discrete Fourier transform (DFT) and its inverse, which can be used for polynomial multiplication.
    \item \textbf{Lagrange Inversion Formula: } A formula that provides the coefficients of the inverse of a power series, which can be applied to polynomials.
    \item \textbf{Lagrange Interpolation: } A method for finding a polynomial that passes through a given set of points.
    \item \textbf{Newton's Method: } An iterative method for finding successively better approximations to the roots of a real-valued function, which can be applied to polynomials.
\end{itemize}
